% Part: sets-functions-relations
% Chapter: sets
% Section: basics

\documentclass[../../../include/open-logic-section]{subfiles}

\begin{document}

\olfileid{sfr}{set}{bas}
\olsection{Extensionaliteit}

Een \emph{verzameling} is een geheel van objecten dat als één
object wordt beschouwd. De objecten waaruit de verzameling bestaat,
heten \emph{elementen} of \emph{leden} van de verzameling. Als $x$
!!a{element} van een verzameling~$A$ is, schrijven we $x \in A$; zo
niet, dan schrijven we $x \notin A$. De verzameling zonder
!!{element}s heet de \emph{lege} verzameling en wordt
aangeduid met~‘$\emptyset$’.

\begin{explain}
Het maakt niet uit hoe we de verzameling \emph{aangeven}, in welke
\emph{volgorde} we haar !!{element}s plaatsen of zelfs hoe
\emph{vaak} we haar !!{element}s meetellen. Alleen welke
!!{element}s ze bevat, is van belang. Dit leggen we vast in het
volgende beginsel.
\end{explain}

\begin{defn}[Extensionaliteit]
  Als $A$ en $B$ verzamelingen zijn, dan geldt $A = B$ dan en slechts
  dan als elk !!{element} van~$A$ ook !!a{element} van~$B$ is, en
  omgekeerd.
\end{defn}

Dankzij extensionaliteit kunnen we bepaalde notatie gebruiken. Zijn
$a_{1}$, \dots, $a_{n}$ objecten, dan is $\{a_{1}, \dots, a_{n}\}$ in
het algemeen \emph{de} verzameling waarvan $a_1, \ldots, a_n$ de
!!{element}s zijn. We benadrukken het woord ‘\emph{de}’, want uit
extensionaliteit volgt dat er maar \emph{één} zo'n verzameling kan
zijn. Uit extensionaliteit volgt bovendien:
  \[
    \{a, a, b\} = \{a, b\} = \{b,a\}.
  \] 
Dit weerspiegelt dat bij verzamelingen de volgorde van de
!!{element}s er niet toe doet, evenmin als het aantal keren dat ze
worden vermeld.

\begin{tagblock}{novice}
\begin{ex}
Elke groep objecten kan tot een verzameling worden samengenomen. De
verzameling van Richards broers en zussen bevat bijvoorbeeld één
persoon en kan worden geschreven als $S=\{\textrm{Ruth}\}$. De
verzameling van de positieve gehele getallen kleiner dan $4$ is
$\{1, 2, 3\}$, maar kan ook worden geschreven als $\{3, 2, 1\}$ of
zelfs als $\{1, 2, 1, 2, 3\}$. Volgens extensionaliteit zijn dit
allemaal dezelfde verzameling: elk !!{element} van $\{1, 2, 3\}$ is
ook !!a{element} van $\{3, 2, 1\}$ (en van $\{1, 2, 1, 2, 3\}$), en
omgekeerd.
\end{ex} 
\end{tagblock}

Vaak geven we een verzameling aan met een eigenschap die haar
!!{element}s gemeen hebben. Daarvoor gebruiken we de volgende verkorte
notatie: $\Setabs{x}{\phi(x)}$. Hier staat $\phi(x)$ voor de
eigenschap die~$x$ moet hebben om tot de !!{element}s van de
verzameling te worden gerekend.

\begin{tagblock}{novice}
\begin{ex}
In ons voorbeeld hadden we $S$ ook als volgt kunnen aangeven:
\[
S = \Setabs{x}{x \text{ een broer of zus van Richard is}}.
\]
\end{ex}
\end{tagblock}

\begin{tagblock}{math}
\begin{ex}
Een getal heet \emph{volmaakt} dan en slechts dan als het gelijk is
aan de som van zijn echte delers (d.w.z. getallen waardoor het zonder
rest deelbaar is, maar die niet gelijk zijn aan het getal zelf). Zo is
$6$ volmaakt, want zijn echte delers zijn $1$, $2$ en~$3$, en
$6 = 1 + 2 + 3$. Het getal $6$ is zelfs het enige positieve gehele
getal kleiner dan $10$ dat volmaakt is. Met extensionaliteit kunnen we
dus schrijven:
\[
	\{6\} = \Setabs{x}{x\text{ volmaakt is en }0 \leq x \leq 10}
\]
We lezen de notatie rechts als ‘de verzameling van alle $x$ waarvoor
$x$ volmaakt is en $0 \leq x \leq 10$’. Deze gelijkheid bevestigt dat
het bij verzamelingen niet uitmaakt hoe ze worden aangegeven. Meer in
het algemeen volgt uit extensionaliteit dat er hoogstens één verzameling
van alle $x$ kan zijn waarvoor $\phi(x)$ geldt. Daarom mogen we
$\Setabs{x}{\phi(x)}$ \emph{de} verzameling van alle $x$ noemen
waarvoor~$\phi(x)$ geldt.
\end{ex}
\end{tagblock}

Met extensionaliteit kunnen we aantonen dat verzamelingen gelijk zijn: om
$A = B$ te bewijzen, tonen we aan dat uit $x \in A$ steeds $x \in B$ volgt
en uit $y \in B$ steeds $y \in A$.

\begin{prob}
Bewijs dat er hoogstens één lege verzameling is, d.w.z. toon aan dat
als $A$ en $B$ verzamelingen zonder !!{element}s zijn, dan $A = B$.
\end{prob}
\end{document}
